%document class
\documentclass[10pt,oneside]{book}
%%%% Page Info + Commands %%%%%{

%packages
\usepackage{pgfplots}
\pgfplotsset{compat=1.18}
\usepackage{amsmath}
\usepackage{amsfonts}
\usepackage{amsthm,letltxmacro}
\usepackage{amssymb}
\usepackage{enumerate}
\usepackage{enumitem}
\usepackage[dvipsnames]{xcolor}
\usepackage{changepage}

%This will shrink the page
\newcommand{\smallpage}{  \setlength \oddsidemargin{.25in}
            \setlength \textwidth{5.5in}
            \setlength \topmargin{0in}
            \setlength \textheight{9in}}


% Commands for sets of numbers
\newcommand{\Z}{\mathbb{Z}}\newcommand{\R}{\mathbb{R}}\newcommand{\N}{\mathbb{N}}

% gordie spacing and boxing commands
\newcommand{\gspc}{\vspace{0.13cm}} % 0.5 space
\newcommand{\gline}{\gspc\\} % 1.5 space
\newcommand{\gbline}{\gspc\gspc\\} % double space
\newcommand{\gbox}[1]{\fbox{\parbox{\linewidth}{#1}}} % multiline box command
\newcommand{\gmod}[1]{\,(\text{mod}\,#1)}


\newcommand{\gimage}[3]{
\begin{figure}[h]   
    \centering          
    \includegraphics[width=#2\textwidth]{#1}  
    \caption{#3}
    \label{fig:#1}
\end{figure}\\
}

\newcommand{\centerit}[1]{{\begin{center} #1 \end{center}}}
\newcommand{\ul}[1]{\underline{#1}}

%%%%%%%% command for graphics %%%%%%%%%%%%%
\usepackage{fancyhdr}
%}

%%%% Page Setup %%%%%%%
\smallpage\sloppy
\pagestyle{fancy}
\lhead{\large{\textbf{Problem Set 8}}}
%\rfoot{\thepage/\pageref{LastPage} }
\setlength{\headheight}{10pt}


% Proof writing
\LetLtxMacro\oldproof\proof
\let\endoldproof\endproof
\renewenvironment{proof}[1][\proofname]
  {\vspace{0.2cm}\begin{adjustwidth}{2em}{2em}
   \oldproof[#1]}
  {\endoldproof
\end{adjustwidth}}

\begin{document}

\newcommand{\cg}[1]{\color{OliveGreen}#1\color{white}}
\newcommand{\cb}[1]{\color{BlueViolet}#1\color{white}}

\subsection*{Last Time}
\begin{itemize}
	\item FLT: If $p$ is prime and $p\nmid a$ then $a^{p-1}\equiv 1\gmod p$. 
	\item Corollary to FLT: If $p$ is prime, then $a^p\equiv a\gmod p$. 
\end{itemize}
Old news: $ax\equiv 1\gmod n$ has a solution if and only if $\gcd(a,n)= 1$. \\
So... if $n=p$ is prime, and $a = 1,2,\ldots, p-1$, $\gcd(a,n)$ is definitely 1. \\
Thus, every $a = 1,2,\ldots, p-1$ ``has an inverse''$\gmod p$.
\subsection*{Inverses}
We begin our discussion of inverses. A key definition is that:\gspc
\begin{adjustwidth}{2em}{2em}
	Two integers $a,b$ are invereses if $ab\equiv 1\gmod p$.
\end{adjustwidth}\gspc
Thus, we have that in $\gmod p$ for prime $p$, every $a \;(1\le a \le p - 1)$ has an inverse $\gmod p$.\gline
Couple of key points:
\begin{itemize}[itemsep=0pt, parsep=0pt]
	\item $a = 1$ and $a = p-1$ are self-inverses. 
	\item What about $a = 2,3,\ldots, p-2$?
\end{itemize}
Claim: $a = 2, 3, \ldots, p-2$ are not self-inverses.
\begin{proof}
	For the sake of contradiction, let $a^2\equiv 1\gmod p$. Thus, we have that $a^2 -1 \equiv 0 \gmod p$. \gline
	It follows that $(a-1)(a+1)=a^2 - 1\equiv 0 \gmod p$\gline
	So, $p\mid (a-1)(a+1)$ so either $p\mid a - 1$ or $p\mid a+1$\gline
	If $2\le a\le p -2$, that implies that $1 \le a- 1 \le p-3$ and $3 \le a+1\le p-1$, but clearly because $p$ is prime $p\nmid a -1$ and $p\nmid a+1$, a contradiction.
\end{proof}
Here's a clear example, $p = 11$:
\begin{align*}
	2\cdot 3\cdot 4\cdot 4 \cdot 5 \cdot 6\cdot 8 \cdot 8 \cdot 9 \cdot 10 \equiv \gmod {11}
\end{align*}
We have that each number has a lil buddy that it likes:
\begin{align*}
	2\cdot 6 &\equiv 1 \mod 11\\
	3\cdot 4 &\equiv 1 \mod 11\\
	5\cdot 9 &\equiv 1 \mod 11\\
	7\cdot 9 &\equiv 1 \mod 11
\end{align*}
This is a neat result, and we have a theorem name for it. 
\subsection*{Wilson's Theorem}
\begin{adjustwidth}{2em}{2em}
	Let $p$ be a prime:
	\begin{gather*}
		(p-1)! \equiv -1 \gmod p
	\end{gather*}
	That's it. That's the theorem. 
\end{adjustwidth}

\end{document}
